\section{Discussion and Limitations}
\label{sec:discussion}

\begin{itemize}
\item Main takeaway: who we fine-tune on changes what Centaur predicts. The model may learn a general answer tendency without getting more accurate or capturing human differences on a new task.
\item The response adjustment helps show what fine-tuning adds beyond making certain answers more likely overall.
\item Differences between adapters can carry over to a new task even when they don't beat frozen Centaur in every condition or match human differences. H2 shows that the population matters; H3 and H4 ask what the model has picked up from it.
\item No consistent Western advantage. Which population is predicted better changes from task to task.
\end{itemize}

\subsection{Limits}
\begin{itemize}
\item Country, language, recruitment, and other sample differences are hard to separate. Nationality is only a rough guide to culture.
\item Few items, small lottery test samples, and only three seeds. Repeated answers don't add more independent people.
\item Better scores can partly mean that the human responses are easier to predict.
\item Comparing with frozen Llama doesn't tell us whether Centaur's original fine-tuning caused its sensitivity to the new fine-tuning population.
\item The datasets were chosen as the project developed. Prior model exposure is also uncertain, so it's unclear how far these findings extend.
\item Giving to Taking is only one test on a new task, and the human difference changes across waves. This doesn't show that transfer fails in general.
\end{itemize}
